\documentclass[10pt,a4paper]{amsart}
\usepackage[margin=19mm]{geometry}
\usepackage{amsmath,amssymb,mathtools}
\usepackage[scr=rsfs,cal=euler]{mathalfa}
\usepackage{xfrac}
\usepackage[unicode,hidelinks,bookmarks=false]{hyperref}
\newcommand{\ie}{i.e.\ }
\newcommand{\cf}{cf.\ }
\newcommand{\resp}{respectively\ }
\newcommand{\Subsection}{\subsection}
\providecommand{\nobreakdash}{\nobreak\hbox{-}}
\setlength{\emergencystretch}{2em}
\multlinegap0pt
\def\A{{\mathcal{A}}}
\usepackage{amssymb}
\usepackage{mathtools}
\usepackage[shortlabels]{enumitem}
\usepackage{caption}
\newtheorem{defn}{Definition}
\newtheorem{theorem}{Theorem}
\DeclareMathOperator{\hochcx}{{HC}}
\DeclareMathOperator{\id}{Id}
\title{The Hochschild homology of a noncommutative symmetric quotient stack}
\author{Rina Anno, Vladimir Baranovsky, Timothy Logvinenko}
\date{Source: arXiv:2512.25039v1 (CC BY 4.0)}
\begin{document}
\maketitle

\noindent\emph{Source and licence.} The Hochschild homology of a noncommutative symmetric quotient stack. Version arXiv:2512.25039v1 (CC BY 4.0), distributed by arXiv under the Creative Commons Attribution 4.0 International licence, http://creativecommons.org/licenses/by/4.0/, as recorded in the arXiv licence metadata for this exact version and shown on the abstract page https://arxiv.org/abs/2512.25039v1. Full licence notice: \texttt{licenses/arxiv-2512.25039-CC-BY-4.0.txt}.\par\smallskip

\noindent\textit{Editorial scope.} Counted unit: the article's theorem theorem-phi-is-the-homotopy-of-gf-minus-id (source0.tex lines 2083--2089). The article's complete original proof of it is delivered with it (source0.tex lines 2090--2431, one \texttt{proof} environment of 342 lines). No mathematics is written, repaired or re-derived: the statement and the proof are the article's own lines, in the article's own notation, and no retained line is substituted or re-flowed.  Retained uncounted as the source-local context the proof needs: lines 1717--1776; lines 1890--2032; lines 2059--2069.  Nothing was omitted from any retained span, so the package carries no omission record.  No retained line contains a citation, so the wrapper carries no bibliography and no bibliography entry.  No retained line contains a figure environment, an image-inclusion command or drawing code, so no image is delivered and the package declares no asset dependency.  Editorial formatting only, in addition: an AMS \texttt{amsart} preamble that restates the article's own declarations and macros used by the retained lines, namely the environments \texttt{defn}, \texttt{theorem} on separate counters, together with the standard packages those lines need.  The retained environments print in this document as Definition 1, Theorem 1 in order of appearance, whereas the article prints the same items with the numbering of its own full document.  The complete native source of the article as distributed in the arXiv e-print for this exact version is bundled unchanged as \texttt{sources/191962/source0.tex}.
\begin{defn}
\label{defn-the-map-g}
Let $\A$ be a small DG category. For each $n \geq 1$ define the map 
$$ g \colon 
\hochcx_\bullet(\A)
\rightarrow 
\hochcx_\bullet(\A^{\otimes n};t)_t
$$
by setting for each $m \geq 1$ the map 
$g \colon 
\hochcx_{m-1}(\A)
\rightarrow 
\hochcx_{m-1}(\A^{\otimes n};t)_t$
to be
\begin{equation}
\label{eqn-quasi-iso-g-long-cycle}
\begin{pmatrix}
\alpha_{1} \\
\alpha_{2} \\
\dots
\\
\alpha_{m} 
\end{pmatrix}
\quad \mapsto \quad 
\sum_{c \in \left\{1,\dots,n\right\}^{m}
\\
\text{ with } c_1 = 1
}
(-1)^{\sigma_{c}}
\begin{pmatrix}
\beta_{11} & \beta_{12} & \dots & \beta_{1n}  
\\
\beta_{21} & \beta_{22} & \dots & \beta_{2n}  
\\
& & \dots & 
\\
\beta_{m1} & \beta_{m2} & \dots & \beta_{mn}.
\end{pmatrix}
\end{equation}
Here: 
\begin{itemize}
\item The summation is taken over all $c = (c_1, \dots, c_m)$ with 
$c_i \in \left\{1, \dots, n\right\}$ and $c_1 = 1$.
We think of each $c_i$ as a choice of a column position in the $i$-th row 
of the matrix $(\beta_{ij})$ and we always choose 
the first position in the first row.
\item We define $\sigma_c \in S_m$ to be 
the permutation where $\sigma_c(i)$ is 
the number of $c_k$ with $c_k < c_i$ or with
$c_k = c_i$ and $k \leq i$. 
\item We define
\begin{equation*}
\beta_{ij} := 
\begin{cases}
\alpha_{\sigma_c(i)} \quad \quad j = c_i, \\ 
\id \quad \quad \; j \neq c_i, \\ 
\end{cases}
\end{equation*}
\end{itemize}
\end{defn}

\begin{defn}
\label{defn-cut-and-slice-matrix-operation-bpqdrc}
Let $m,n \geq 1$. We define the map 
\begin{equation}
B_{p,q,d,r,\underline{c}}\colon 
\hochcx_{m-1}(\A^{\otimes n})  \rightarrow 
\hochcx_{m}(\A^{\otimes n})
\end{equation}
parametrised by
\begin{itemize}
\item $p \in \left\{1, \dots, m\right\}$,
\item $q \in \left\{2, \dots, n\right\}$,
\item $d \in \left\{1, \dots, p\right\}$, 
\item $r \in \left\{0, \dots, d-1\right\}$ if $q < n$ and $r = 0$ if $q = n$, 
\item $\underline{c} = (c_1, \dots, c_r) \in \left\{ 1, \dots, n-q \right\}^r$,
\end{itemize}
to be the map 
which sends any chain
\begin{equation*}
\underline{\alpha}
:= 
\begin{pmatrix}
\alpha_{11} & \alpha_{12} & \dots & \alpha_{1n}  
\\
\alpha_{21} & \alpha_{22} & \dots & \alpha_{2n}  
\\
& & \dots & 
\\
\alpha_{m1} & \alpha_{m2} & \dots & \alpha_{mn}. 
\end{pmatrix} 
\end{equation*}
to the chain $\underline{\beta} = (\beta_{ij})$ defined by 
\\
\begin{tabular}{|l|l|l|l|}
\hline
Row group:
&
$i$
&
$j$
&
$\beta_{ij}$
\\
\hline
$1$
&
$1$ 
&
$1$
&
$\alpha_{(p+1)q}\dots\alpha_{mn}\alpha_{11}$
\\
&
&
otherwise
&
$\id$
\\
\hline
$2$
&
$2,\dots,r+1$
&
$c_{i-1}$
&
$\alpha_{\left(\sigma_{\underline{c}}(i-1)+1\right)1}$
\\
&
&
otherwise
&
$\id$
\\
\hline
$3$
&
$r+2, \dots, d$
& 
$n-q+1$
&
$\alpha_{i1}$
\\
& 
&
otherwise
&
$\id$ 
\\
\hline
$4$
&
$d+1, \dots, m+d-p$
&
$n-q+2,\dots, n$
&
$\alpha_{(i+p-d)(j-n+q-1)}$
\\
&
&
otherwise
& 
$\id$  
\\
\hline
$5$
&
$(m+d-p+1)$
& 
$n-q+2,\dots,n$
&
$\alpha_{1(j-n-q)}\dots\alpha_{d(j-n-q)}$  
\\
&
&
otherwise
&
$\id$  
\\
\hline
$6$
&
$m+d-p+2, \dots, m+1$
& 
$n-q+1,\dots,n$  
&
$\alpha_{(i-m+p-1)(j-n+q)}$ 
\\
&
& 
otherwise
&
$\id$ 
\\
\hline
\end{tabular}
\\
where $\sigma_{\underline{c}} \in S_r$, similar to 
Defn.~\ref{defn-the-map-g}, is the permutation of the rows
which arises when we distribute  
$\alpha_{21}, \dots, \alpha_{(r+1)1}$ in composable order
onto the positions in the rows $2,\dots,r+1$ chosen by
$\sigma_{\underline{c}}$. 
\end{defn}

\begin{equation}
\label{eqn-defn-of-psi'-n}
\sum_{p, q, d, r, \underline{c}}
\sum_{\tau }
\sum_{\upsilon}
(-1)^{\sigma_{\underline{c}}}
(-1)^{\tau}
(-1)^{\upsilon} 
(-1)^{m+(m+d)(p+d)}
\tau \upsilon B_{p,q,d,r,\underline{c}} \left( \underline{\alpha} \right). 
\end{equation}

\begin{theorem}
\label{theorem-phi-is-the-homotopy-of-gf-minus-id}
We have an equality of maps 
$\hochcx_\bullet(\A^{\otimes n}; t)_t 
\rightarrow \hochcx_\bullet(\A^{\otimes n}; t)_t$:
$$ d\Phi_n = g_n f_n - \id. $$
\end{theorem}

\begin{proof}
Let $f'_n$ be the $s = 1$ summand of 
$f_n$, so that $f_n = \frac{1}{n} \sum_{s=1}^n f'_n \circ t^{s-1}$. 
It suffices to show that we have an equality of maps 
$\hochcx_\bullet(\A^{\otimes n}; t) 
\rightarrow \hochcx_\bullet(\A^{\otimes n}; t)$:
$$ d\Phi'_n = g_n f'_n - \id. $$

For any $m \geq 1$ take any basic Hochschild chain 
$\alpha \in \hochcx_{m-1}(\A^{\otimes n})$ 
\begin{equation*}
\underline{\alpha}
:= 
\begin{pmatrix}
\alpha_{11} & \alpha_{12} & \dots & \alpha_{1n}  
\\
\alpha_{21} & \alpha_{22} & \dots & \alpha_{2n}  
\\
& & \dots & 
\\
\alpha_{m1} & \alpha_{m2} & \dots & \alpha_{mn}. 
\end{pmatrix} 
\end{equation*}

We need to show that 
\begin{equation}
\label{eqn-dpsi-equals-minus-psid+gf-id}
d\left(\Phi'_n\left(\underline{\alpha}\right)\right) = 
- \Phi'_n\left(d\underline{\alpha}\right) + g_nf'_n(\underline{\alpha}) - 
\underline{\alpha}. 
\end{equation}

We verify this by going through each summand in 
the sum \eqref{eqn-defn-of-psi'-n} for $\Phi'_n(\underline{\alpha})$ on the LHS
of \eqref{eqn-dpsi-equals-minus-psid+gf-id}
and differentiating it. Each summand of the result 
is a basic chain of $\hochcx_{m-1}(\A^{\otimes n})$ given 
by an $(m \times n)$-matrix involving $\alpha_{ij}$ 
with a coefficient $\pm 1$ coming from the the sign twists in 
\eqref{eqn-defn-of-psi'-n} and in the Hochschild differential. 
For each such matrix, we point
out that it occurs exactly one more time in the
sum \eqref{eqn-defn-of-psi'-n} for $d(\Phi'_n(\underline{\alpha}))$ 
on the LHS, in the sum \eqref{eqn-defn-of-psi'-n} for $\Phi'_n(d
\underline{\alpha})$ on the RHS, 
or in the term $g_nf'_n(\underline{\alpha}) - \underline{\alpha}$. 
We point out where it occurs and verify that the signs are such 
that the cancellation occurs. 

At the end of the proof, we will list the special summands of 
$d(\Phi'_n(\underline{\alpha}))$ on the LHS which cancel out 
those in the the term  $g_nf'_n(\underline{\alpha}) -
\underline{\alpha}$. For most of the proof, we will take 
the summands of $d(\Phi'_n(\underline{\alpha}))$ on the LHS, and 
identify the corresponding summands of 
$d(\Phi'_n(\underline{\alpha}))$ on the LHS or of 
$\Phi'_n(d \underline{\alpha})$ on the RHS which cancel them out.  
Our analysis makes it clear that all of the summands 
of $d(\Phi'_n(\underline{\alpha}))$ on the RHS 
are cancelled out. 

We first assume that $\tau = \upsilon = \id_{S_{m+1}}$. This is the
main substance of our analysis. We now consider the following cases: 

\fbox{$p, q, d, r$, and $\underline{c}$ generic:}

We start with generic values of $p, q, d, r$, and $\underline{c}$
in \eqref{eqn-defn-of-psi'-n} for which all six row
ranges listed in the table definining $B_{p, q, d, r
\underline{c}}(\underline{\alpha})$ in
Defn.~\ref{defn-cut-and-slice-matrix-operation-bpqdrc} are present and
non-empty. Explicitly, this means that $r \geq 2$, $d \geq r+1$, 
$d < p < m$, and $q < n$. 

We produce the cancelling summand for each summand of
$d B_{p, q, d, r \underline{c}}(\underline{\alpha})$. 
Denote by $d_i B_{p, q, d, r \underline{c}}(\underline{\alpha})$
the summand where $i$th and $(i+1)$st rows of
$B_{p, q, d, r \underline{c}}(\underline{\alpha})$ are merged.  
The way we produce the cancelling summand for 
$d_i B_{p, q, d, r \underline{c}}(\underline{\alpha})$
depends only on the row group $i$ and the row group of $i+1$ 
in the table of the row groups of 
$B_{p, q, d, r \underline{c}}(\underline{\alpha})$
given in Definition
\ref{defn-cut-and-slice-matrix-operation-bpqdrc}. 
It thus suffices to check the following cases, where by
case $(pq)$ we mean the case where $i$ lies in group $p$
and $i+1$ lies in group $q$:
\begin{itemize}
\item \underline{(12)}: This is the case $i = 1$. It has  
two subcases. If $c_1 = 1$, the 
the cancelling summand is  
$$B_{p-1, q, d-1, r-1,(c_2, \dots, c_r)}(d_1 \underline{\alpha}).$$
If $c_1 > 1$, then it is 
$$
d_{m+1} \tau B_{p, q, d, r, (c_2, \dots, c_r, c_1 - 1)}(\underline{\alpha})
\quad \text{ with } \quad \tau = \left((m+1)m\dots(r+1)\right).$$

\item \underline{(22)}: There are two subcases. If $c_i =
c_{i+1}$, the cancelling summand is 
$$B_{p-1, q, d-1, r-1,(c_1, \dots, c_{i}, c_{i+2}, \dots,
c_r)}(d_{\sigma_{\underline{c}}(i)} \underline{\alpha}).$$ 
If $c_i \neq c_{i+1}$, then it is 
$$d_{i} B_{p, q, d, r,(c_1, \dots, c_{i-1},
c_{i+1}, c_{i}, c_{i+2} \dots, c_r)} (\underline{\alpha}).$$ 

\item \underline{(23)}: This is the case $i = r+1$. 
The cancelling summand is 
$$ d_{r+1} \tau B_{p, q, d, r,\underline{c}} (\underline{\alpha})
\quad \text{ with } \quad \tau = \left((r+1)(r+2)\right). $$

\item \underline{(33)}: The cancelling summand is 
$$ B_{p-1, q, d-1, r, \underline{c}} (d_i \underline{\alpha}). $$

\item \underline{(34)}: This is the case $i=d$ The cancelling summand is 
$$ d_d \upsilon B_{p, q, d, r, \underline{c}} (\underline{\alpha})
\quad \text{ with } \quad \upsilon = \left(d(d+1)\right). $$

\item \underline{(44)}: The cancelling summand is 
$$ B_{p, q, d, r, \underline{c}} (d_{i + p - d} \underline{\alpha}).$$

\item \underline{(45)}: This is the case $i = m-p+d$. 
The cancelling summand is 
$$ B_{p, q, d, r, \underline{c}} (d_{m} \underline{\alpha}).$$

\item \underline{(56)}: This is the case $i = m-p+d+1$. 
The cancelling summand is 
$$ d_{m-p+d+1} \upsilon B_{p, q, d+1, r, \underline{c}} (\underline{\alpha})
\quad \text{ with } \quad \upsilon = \left((m-p+d+1)\dots(d+1)\right). $$

\item \underline{(66)}: The cancelling summand is 
$$ B_{p-1, q, d, r, \underline{c}} (d_{i-m+p-1} \underline{\alpha}).$$

\item \underline{(61)}: This is the case $i = m+1$. 
The cancelling summand is 
$$ d_1 \tau \upsilon B_{p-1, q, d, r, \underline{c}} (\underline{\alpha})
\quad \text{ with } \quad 
\tau = \left(2\dots(r+2)\right),
\upsilon = \left((r+2)\dots(d+1)\right). $$
\end{itemize}

\fbox{$p, q, d, r$, and $\underline{c}$ non generic:}

These are the cases when one or more of the six row
ranges listed in the table definining $B_{p, q, d, r
\underline{c}}(\underline{\alpha})$ in
Defn.~\ref{defn-cut-and-slice-matrix-operation-bpqdrc} are not
present. 

This adds to the above analysis the following additional cases 
where we need to describe the cancelling summand: 

\begin{itemize}
\item \underline{(13)}: This is the case $r = 0$ but $d \geq 2$. 
There are two subcases. If $q = n$, the the cancelling summand is  
$$B_{p-1, n, d-1, 0,\emptyset}(d_1 \underline{\alpha}),$$
where $\emptyset$ is the value of $\underline{c}$ when $r=0$. 
If $q < n$, the cancelling summand is
$$
d_{m+1} \tau B_{p, q, d, 1, (n-q)}(\underline{\alpha})
\quad \text{ with } \quad \tau = \left((m+1)m\dots2\right).$$

\item \underline{(14)}: This is the case $r = 0$, $d = 1$, $p < m$. The
cancelling summand is 
$$ B_{p+1, q, 1, 0, \emptyset}(d_{m+1}\underline{\alpha}). $$

\item \underline{(15)}: This is the case $r = 0$, $d = 1$, $p = m$ 
The are two cases. If $q = n$, then the cancelling summand is 
$\underline{\alpha}$ itself on the RHS of \eqref{eqn-defn-of-psi'-n}. 
If $q \neq n$, then the cancelling summand is
$$ B_{1, q+1, 1, 0, \emptyset}(d_{m+1}\underline{\alpha}). $$

\item \underline{(24)}: This is the case $d > 1$, $r = d-1$, $p < m$.
The cancelling summand is 
$$ d_d \tau B_{p, q, d, d-1, \underline{c}} (\underline{\alpha})
\quad \text{ with } \quad 
\tau = \left(d(d+1)\right).$$

\item \underline{(25)}: This is the case $d > 1$, $r = d-1$, $p = m$. 
The cancelling summand is 
$$ d_{d} \tau B_{m, q, d, d-1, \underline{c}} (\underline{\alpha})
\quad \text{ with } \quad 
\tau = \left(d(d+1)\right).$$

\item \underline{(35)}: This is the case $d > 1$, $r < d-1$, $p = m$.
The cancelling summand is:
$$ d_{d} B_{m,q,d-1,r,\underline{c}}(\underline{\alpha}). $$

\item \underline{(51)}: This is the case $d=p$. There are two
subcases. If $q = 2$, the cancelling summand is 
the summand of $g_n f'_n(\underline{\alpha})$ corresponding
to the summand of $g_n$ indexed by 
\begin{equation}
\label{eqn-indexing-c-for-the-summand-of-gf'}
\underline{c}' = (1,c_1, \dots, c_r,n-1, \dots, n-1, n, \dots, n) \in 
\left\{ 1, \dots, n\right\}^m,
\end{equation}
where there are $p-1-r$ and $m-p$ instances of $n-1$ and $n$,
respectively. If $q \neq 2$, the cancelling summand is
$$ d_{1} \tau B_{m,q-1,d,r,(c_1, \dots, c_r, n-q+1, \dots, n-q+1)}(\alpha)
\quad \text{ with } \quad \tau = (23\dots(d+1)).$$
\end{itemize}

Finally, let us assume that $\upsilon$ is non-trivial. 
Then 
$d \left(\upsilon B_{p,q,d,r,\underline{c}}(\underline{\alpha})\right)$ 
has three kinds of summands. First are those where the differential 
merges two rows belonging one each to the two row subsets which
$\upsilon$ shuffles. Then let $\upsilon'$ be the composition of
$\upsilon$ with the transposition interchanging these two rows. Then 
$d \left(\upsilon' B_{p,q,d,r,\underline{c}}(\underline{\alpha})\right)$ 
has an identical summand coming from merging the same two rows. All 
the sign twists in \eqref{eqn-defn-of-psi'-n} stay the same, 
except for $(-1)^{\upsilon'} = - (-1)^{\upsilon}$. We conclude that
the two summands cancel each other out. 

The second type of summands in 
$d \left(\upsilon B_{p,q,d,r,\underline{c}}(\underline{\alpha})\right)$
are those which do not belong to the first group, but
where the two rows being merged were also adjacent in 
$B_{p,q,d,r,\underline{c}}(\underline{\alpha})$. 
Thus there is a summand of 
$d(B_{p,q,d,r,\underline{c}}(\underline{\alpha}))$ 
where these two rows are merged.  
The analysis above give the cancelling counterpart $X$ for this summand. 
Let $\upsilon' \in S_m$ be the permutation obtained from $\upsilon$
where we treat the two merged rows as a single row. 
Then $\upsilon' X$ would cancel the desired summand of 
$d \left(\upsilon B_{p,q,d,r,\underline{c}}(\underline{\alpha})\right)$. 
It remains to show that $\upsilon'X$ exists as a summand in  
\eqref{eqn-dpsi-equals-minus-psid+gf-id}. We verify this
by analysing every applicable cancelling summand in the analysis
above for the case $\tau = \upsilon = \id$. Firstly, when 
$X = B_\bullet(d_\bullet \underline{\alpha})$, it is clear 
that $\upsilon' X$ is also a summand of $\Phi'_n(d \alpha)$. 
The remaining cases are:
\begin{itemize}
\item \underline{(12)}, the subcase $c_1 > 1$. Then 
$$ \upsilon' X = \upsilon' 
d_{m+1} \tau B_{p, q, d, r, (c_2, \dots, c_r, c_1 - 1)}(\underline{\alpha})
= 
d_{m+1} \tau \upsilon 
B_{p, q, d, r, (c_2, \dots, c_r, c_1 - 1)}(\underline{\alpha})
$$
where $\tau = \left((m+1)m\dots(r+1)\right)$.
\item \underline{(22)}, the subcase $c_i \neq c_{i+1}$. Then 
$$
\upsilon' X = \upsilon' 
d_{i} B_{p, q, d, r,(c_1, \dots, c_{i-1},
c_{i+1}, c_{i}, c_{i+2} \dots, c_r)} (\underline{\alpha})
= 
d_{i} \upsilon B_{p, q, d, r,(c_1, \dots, c_{i-1},
c_{i+1}, c_{i}, c_{i+2} \dots, c_r)} (\underline{\alpha}).$$ 
\item \underline{(23)}. Then 
$$ \upsilon' X = 
\upsilon' d_{r+1} \tau B_{p, q, d, r,\underline{c}}(\underline{\alpha})
= 
d_{r+1} \tau \upsilon B_{p, q, d,
r,\underline{c}}(\underline{\alpha}), $$
where $\tau = \left((r+1)(r+2)\right)$. 
\item \underline{(56)}. Then 
$$ \upsilon' X = 
\upsilon' d_{m-p+d+1} \upsilon'' B_{p, q, d+1, r, \underline{c}}
(\underline{\alpha})
= d_{m-p+d+1} \upsilon \upsilon'' B_{p, q, d+1, r, \underline{c}}
(\underline{\alpha}) $$
where $\upsilon'' = \left((m-p+d+1)\dots(d+1)\right).$
\item \underline{(61)}. Then
$$ \upsilon' X = 
\upsilon' d_1 \tau \upsilon'' B_{p-1, q, d, r, \underline{c}} (\underline{\alpha})
=
d_1 \tau (t_{m+1} \upsilon t_{m+1}^{-1} \upsilon'')  B_{p-1, q, d, r, \underline{c}} (\underline{\alpha})
$$
where $\tau = \left(2\dots(r+2)\right), 
\upsilon'' = \left((r+2)\dots(d+1)\right), t_{m+1} = (1\dots{m+1}).$
\item \underline{(13)}, the subcase $q < n$.
Then 
$$ \upsilon' X =
\upsilon'
d_{m+1} \tau B_{p, q, d, 1, (n-q)}(\underline{\alpha})
= d_{m+1} \upsilon \tau B_{p, q, d, 1, (n-q)}(\underline{\alpha}), $$
where $\tau = \left((m+1)m\dots2\right)$.

\item \underline{(51)}. In the subcase $q = 2$, recall that
$X$ is the summand of $g_n f'_n(\underline{\alpha})$ 
corresponding to the summand of $g_n$ indexed by $\underline{c}'$
as in \eqref{eqn-indexing-c-for-the-summand-of-gf'}. 
Then $\upsilon' X$ is the summand of 
$g_n f'_n(\underline{\alpha})$ corresponding to $\upsilon' \underline{c}'$. 

In the subcase $q \neq 2$, 
\begin{align*}
\upsilon' X & = 
\upsilon' d_{1} \tau B_{m,q-1,d,r,(c_1, \dots, c_r, n-q+1, \dots,
n-q+1)}(\underline{\alpha}) = \\
& = 
d_{1} \tau' B_{m,q-1,d,r,(c_1, \dots, c_r, n-q+1, \dots,
n-q+1)}(\underline{\alpha}),
\end{align*}
where $\tau = (23\dots(d+1))$ and $\tau' = t_{m+1} \upsilon
t_{m+1}^{-1} \tau$. 
\end{itemize}
 
The third and the last type of summands are those which do not belong
to the first group and where the two rows being merged were not adjacent in 
$B_{p,q,d,r,\underline{c}}(\underline{\alpha})$. This necessarily
means that one of the rows being merged is involved in the shuffle $\upsilon$, 
and the other isn't. There are three such cases, where by the case
$(ij)$ we mean the case where the rows being merged, before
they were permuted by $\upsilon$, belonged to the row
groups $i$ and $j$ of the table in 
Definition \ref{defn-cut-and-slice-matrix-operation-bpqdrc}:
\begin{itemize}
\item \underline{(24)}: This occurs when the shuffle $\upsilon$
permutes $d+1$-st row (the top row in group 4) just below the $r+1$-st
row (the bottom row of the group 2). The cancelling summand to
$d_{r+1} \left(\upsilon B_{p,q,d,r,\underline{c}}(\underline{\alpha})\right)$
is then 
$$ d_{r+1} \left(\tau \upsilon
B_{p,q,d,r,\underline{c}}(\underline{\alpha})\right) \quad \quad
\text{ with } \tau = ((r+1)(r+2)).$$
\item \underline{(35)}: This occurs when $\upsilon$ shuffles the
bottom row of the 3rd row group all the way down to the row $m-p+d+1$, 
just above the $5$th row group.  The cancelling summand is 
$$ \upsilon d_{d + m - p} B_{p,q,d-1,r,\underline{c}}(\underline{\alpha}). $$
where $\upsilon' \in S_m$ is obtained from $\upsilon \in S_{m+1}$ by
skipping $d$ in $\{1, \dots, m+1\}$ on the source side, and 
$\upsilon(d)$ on the target side.
\item \underline{(14)}: This occurs when the 2nd row group is empty, 
3rd row group is present, but $\upsilon$ shuffles the top row of the
$4$th group all the way up to the second row. 
The cancelling summand is 
$$ \upsilon' B_{p+1, q, d, r, \emptyset}(d_{m+1}\underline{\alpha}), $$
where $\upsilon' \in S_m$ is obtained from $\upsilon \in S_{m+1}$ by
skipping $d+1$ in $\{1, \dots, m+1\}$ on the source side, and 
$\upsilon(d+1)$ on the target side. 
\end{itemize}

Similar analysis reduces the case of arbitrary $\tau$ to the case 
$\tau = \id$.  
\end{proof}

\end{document}
